Text reaches language models with formatting the user never sees: Windows line endings, no-break spaces copied from web pages, trailing spaces, and typographic quotes and dashes from word processors.
We ask whether such meaning-preserving edits change model outputs, and which kind matters more.
We apply ten invisible whitespace and punctuation edits to three classification tasks and two code-generation tasks.
We run 11 open models from 0.5B to 14.8B parameters with greedy decoding (239,616 generations), and we compare every edit against a byte-identical re-run of the same prompt.
A single invisible edit flips 4.9\% of classification labels (95\% CI 4.6--5.1\%), while the identical re-run flips 0.02\%.
Contrary to our hypothesis, whitespace edits flip 7.1\% of labels and punctuation edits only 2.2\%. Whitespace is more disruptive for every model and on both task types, and it stays more disruptive when we match the number of prompt tokens the edit changes.
Models under 3B parameters flip labels 2.2$\times$ as often as larger ones (CI 2.0--2.4).
For code, 36\% of generations change text under an edit, but only 5.1\% change their unit-test outcome, about the same rate as classification label flips.
Net accuracy moves by less than 1 percentage point per model, because flips go in both directions and concentrate on items the model already gets wrong.
Formatting therefore adds variance rather than bias. Practitioners should normalise whitespace before inference and report results over several formattings.
